\documentclass{beamer}
\usetheme{Madrid}
\usecolortheme{beaver}
\setbeamertemplate{navigation symbols}{}

\usepackage[utf8]{inputenc}
\usepackage[T1]{fontenc}
\usepackage[polish]{babel}
\usepackage{booktabs}
\usepackage{array}
\usepackage{listings}
\usepackage{hyperref}

\title{Co po matematyce?}
\author{Vlad Snozyk}
\date{}

\begin{document}

\begin{frame}
  \titlepage
\end{frame}

\begin{frame}{Wprowadzenie}
  \begin{itemize}
    \item Studia na uczelni
    \item Pierwsza praca
    \item Ogólne doświadczenia zawodowe
  \end{itemize}
\end{frame}

\begin{frame}{Plan prezentacji}
  \begin{itemize}
    \item Preferencje osobiste
    \item Wybory i ścieżka kariery
  \end{itemize}
\end{frame}

\begin{frame}{Ulubione zajęcia}
  \begin{columns}[T]
    \begin{column}{0.33\textwidth}
      \textbf{Świetne}
      \begin{itemize}
        \item Matematyka dyskretna
        \item Teoria prawdopodobieństwa
        \item Sieci neuronowe
        \item Algebra
        \item Topologia
        \item Algebra abstrakcyjna
        \item Sterowanie stochastyczne
        \item Matematyka finansowa
      \end{itemize}
    \end{column}
    \begin{column}{0.33\textwidth}
      \textbf{W porządku}
      \begin{itemize}
        \item Analiza matematyczna
        \item Równania różniczkowe
        \item Programowanie w Pythonie
        \item Programowanie w~R
        \item Logika matematyczna
        \item Statystyka
        \item Procesy stochastyczne
      \end{itemize}
    \end{column}
    \begin{column}{0.33\textwidth}
      \textbf{Nieprzypadły do gustu}
      \begin{itemize}
        \item Programowanie w~Excelu
        \item Ekonometria
        \item Mikroekonomia
        \item Analiza IV
        \item Programowanie w~MATLAB-ie
      \end{itemize}
    \end{column}
  \end{columns}
\end{frame}

\begin{frame}{Przydatność zajęć w pierwszej pracy}
  \begin{columns}[T]
    \begin{column}{0.34\textwidth}
      \textbf{Bardzo przydatne}
      \begin{itemize}
        \item Ekonometria
        \item Statystyka
        \item Programowanie w~R
        \item Programowanie w~Pythonie
        \item Programowanie w~Excelu
        \item Matematyka finansowa
      \end{itemize}
    \end{column}
    \begin{column}{0.33\textwidth}
      \textbf{Trochę przydatne}
      \begin{itemize}
        \item Teoria prawdopodobieństwa
        \item Algebra
        \item Analiza matematyczna
        \item Procesy stochastyczne
      \end{itemize}
    \end{column}
    \begin{column}{0.33\textwidth}
      \textbf{Nieprzydatne}
      \begin{itemize}
        \item Matematyka dyskretna
        \item Topologia
        \item Algebra abstrakcyjna
        \item Sterowanie stochastyczne
        \item Równania różniczkowe
      \end{itemize}
    \end{column}
  \end{columns}
\end{frame}

\begin{frame}{Przydatne zajęcia w obecnej pracy}
  \begin{itemize}
    \item Statystyka
    \item Programowanie w~Pythonie
    \item Teoria prawdopodobieństwa
    \item Algebra
    \item Sterowanie stochastyczne
    \item Matematyka finansowa
    \item Analiza matematyczna
    \item Metody Monte Carlo
  \end{itemize}
\end{frame}

\begin{frame}{Moja ścieżka kariery}
  \begin{itemize}
    \item UBS
    \item Superfund
    \item Citibank
    \item Pentech AI
  \end{itemize}
\end{frame}

\begin{frame}{UBS}
  \begin{itemize}
    \item Proces rekrutacji
      \begin{itemize}
        \item Matematyka finansowa: 20\%
        \item Ekonometria: 40\%
        \item Programowanie w~R: 40\%
      \end{itemize}
    \item Co było dobre?
    \item Co było słabe?
    \item Zabawna historia o języku angielskim
    \item Dlaczego i jak odszedłem?
  \end{itemize}
\end{frame}

\begin{frame}{Pytanie do publiczności}
  Kto polubił historię o języku angielskim?
\end{frame}

\begin{frame}{Superfund}
  \begin{itemize}
    \item Proces rekrutacji
    \item Co było dobre?
      \begin{itemize}
        \item Ciekawostka o programowaniu funkcyjnym
      \end{itemize}
    \item Co było słabe?
    \item Dlaczego i jak odszedłem?
  \end{itemize}
\end{frame}

\begin{frame}{Monoidy: matematyka i programowanie}
  \begin{block}{Definicja własności łączności}
    Działanie $\oplus$ jest łączne, jeśli \((a \oplus b) \oplus c = a \oplus (b \oplus c)\).
  \end{block}
  \begin{block}{Przykład z maksimum}
    \begin{itemize}
      \item $\max(\max(a,b),c) = \max(a,\max(b,c))$
      \item $\max(\max(3,7),5) = \max(7,5) = 7$
      \item $\max(3,\max(7,5)) = \max(3,7) = 7$
    \end{itemize}
  \end{block}
\end{frame}

\begin{frame}{Maksimum listy}
  \begin{block}{Przykład}
\begin{verbatim}
maximum([3, 7, 5, 2, 9, 1]) = 9
\end{verbatim}
  \end{block}
\end{frame}

\begin{frame}{Pytanie do publiczności}
  Kto polubił krótkie wprowadzenie do monoidów?
\end{frame}

\begin{frame}{Citibank}
  \begin{itemize}
    \item Proces rekrutacji
    \item Co było dobre?
    \item Co było słabe?
      \begin{itemize}
        \item Dlaczego było naprawdę źle
        \item Zagadki rekrutacyjne o owocach
      \end{itemize}
    \item Dlaczego i jak odszedłem?
  \end{itemize}
\end{frame}

\begin{frame}{Pytanie do publiczności}
  Kto polubił zagadkę o owocach?
\end{frame}

\begin{frame}{Co rozczarowuje w pracy w banku}
  \begin{itemize}
    \item \url{https://www.youtube.com/shorts/4STft4geMWY}
    \item Jakość modeli
      \begin{itemize}
        \item Cena ropy naftowej
      \end{itemize}
    \item Znaczenie modeli
      \begin{itemize}
        \item Procenty mnożone przez~2
      \end{itemize}
    \item Potencjalne obowiązki
      \begin{itemize}
        \item 240 arkuszy kalkulacyjnych
      \end{itemize}
  \end{itemize}
\end{frame}

\begin{frame}{Pentech AI}
  \begin{itemize}
    \item Jak dostałem tę pracę?
    \item Czym się obecnie zajmuję?
    \item Najlepsze elementy mojej pracy
      \begin{itemize}
        \item Narzędzia
        \item Kultura pracy zdalnej
        \item Problemy do rozwiązania
      \end{itemize}
  \end{itemize}
\end{frame}

\begin{frame}{Pytania}
  \centering
  Czy są pytania?
\end{frame}

\begin{frame}{Dziękuję}
  \centering
  Dziękuję za uwagę!
\end{frame}

\end{document}
